\section{Related Work}

Safety alignment and over-refusal in Large Language Models (LLMs) have become critical areas of research as models are deployed in more diverse settings. Our work builds on several key research directions.

\para{Safety Alignment and Over-refusal} LLMs are typically aligned for safety through Reinforcement Learning from Human Feedback (RLHF) and other preference-based techniques. While effective at reducing harmful outputs, these methods often result in exaggerated safety, where models decline benign requests. \citet{rottger2023xstest} introduced \xstest, a manually curated benchmark designed to identify such over-refusal, particularly for borderline cases. Similarly, \citet{cui2024orbench} released \orbench, a large-scale dataset for evaluating over-refusal by rewriting harmful prompts into benign versions. These benchmarks highlight that current safety-aligned models frequently fail to distinguish between harmful intent and benign context. Unlike these works, which focus on \emph{measuring} over-refusal, we investigate how easily it can be \emph{induced} and its generalization properties.

\para{Mechanistic Understanding of Refusal} Recent studies have sought to understand the internal mechanisms that drive refusal behavior. \citet{lieberum2024refusal} utilized Sparse Autoencoders (SAEs) to identify specific directions in the activation space that correspond to refusal. They found that refusal is a highly steerable behavior, suggesting that it is tied to a well-defined internal representation. This is consistent with our hypothesis that fine-tuning on a small set of benign refusals can shift the entire refusal boundary by amplifying these internal features.

\para{Mitigating Over-refusal} Several methods have been proposed to reduce over-refusal without compromising safety. \citet{dabas2025justenough} introduced ACTOR, a representation fine-tuning technique that shifts the representations of safe queries away from "refusal directions." This indicates that the "refusal boundary" is sensitive to targeted interventions. Our work complements this by showing that this sensitivity is a double-edged sword: while it allows for targeted mitigation, it also enables a rapid and broad induction of over-refusal from very little training data.

\para{Generalization of Behaviors in LLMs} The generalization of behaviors during fine-tuning is a core question in machine learning. Previous work has explored how models learn new tasks or styles from small datasets. Our study focuses specifically on the generalization of \emph{refusal}, which we find to be particularly "contagious." This suggests that safety alignment might operate through a highly generalized mechanism that is prone to over-expansion.
